\chapter{Introduzione}
\label{cap:introduzione}

\section{Obiettivo}

Il progetto realizza un nucleo di calcolo per il prodotto fra una matrice densa
e un multivettore,
\begin{equation}
  Y \leftarrow A X ,
  \qquad
  A \in \mathbb{R}^{M \times N},\quad
  X \in \mathbb{R}^{N \times k},\quad
  Y \in \mathbb{R}^{M \times k},
  \label{eq:problema}
\end{equation}
con $k$ \emph{piccolo}: il multivettore è una matrice densa alta e strettissima.
Il codice funziona per $k$ generico ed è collaudato sui cinque valori richiesti
dalla traccia, $k \in \{3, 6, 8, 20, 32\}$.

\medskip

La parallelizzazione è ibrida a due livelli:

\begin{itemize}
  \item \textbf{MPI (Message Passing Interface)} costituisce il livello esterno, responsabile della distribuzione dei dati su una griglia di processi bidimensionale e della comunicazione fra processi
  \item la parallelizzazione interna al singolo processo è affidata a \textbf{CUDA (Compute Unified Device Architecture)},
  che costituisce il backend di accelerazione su GPU NVIDIA.
\end{itemize}
I due livelli non sono alternativi ma sovrapposti: ogni processo
MPI possiede un blocco di $A$ e delega il calcolo del proprio contributo
parziale a un \emph{backend} locale, che può essere un kernel CUDA oppure un nucleo di CPU.